
\begin{thebibliography}{99}
  \bibitem{ref1} R. J. Baker (2000) CMOS Circuit Design, Layout and Simulation, 3rd edn. PHI, New Delhi
  \bibitem{ref2} Behzad Razavi, (2000) Design of Analog CMOS Integrated Circuits, Tata-McGraw Hill, New Delhi
  \bibitem{ref3} S. A. Loan, M. Nizamuddin, A. R. Alamoud and S.A. Abbasi (2015) "Design and Comparative Analysis of High-Performance Carbon Nanotube-Based Operational Transconductance Amplifiers”, World Scientific Journal Nano, 10(3), 1-11
  \bibitem{ref4} Bano, N., Nizamuddin, M., Ahmad, W., Hashmi, F. and Alattab, A.A., (2025) “Design and performance of graphene nanoribbon FET-based OTA for fast and energy-efficient electronics in IoT” Journal of Nanoelectronics and Optoelectronics, 20(2), pp.181-192.
  \bibitem{ref5} R. L. Geiger and E. Sánchez-Sinencio (1985) “Active Filter Design Using Operational Transconductance Amplifiers: A Tutorial" IEEE Circuits and Devices Magazine, Vol. 1, pp.20-32
  \bibitem{ref6} M. Nizamuddin, S. A. Loan, A. R. Alamoud. and S. A. Abbasi (2015) “Design, simulation and comparative analysis of CNT based Cascode Operational Transconductance Amplifiers”, Nanotechnology, 26 (12pp)
  \bibitem{ref7} Hashmi F., Nizamuddin M., Ahmad W., (2023) “A novel carbon nanotube field effect transistors based   cascode operational transconductance amplifier: An optimum design” Journal of Nanoelectronics and Optoelectronics, 18 (5) 534–546
  \bibitem{ref8} Hashmi, F., Nizamuddin, M. and Amin, S.U., (2025) “Performance evaluation and optimization of triple cascode operational transconductance amplifiers using GNRFET technology for low power smart devices” Materials Research Express, 12(2), p.025601.
  \bibitem{ref9} S. Liu, Z. Zhu, J. Wang, L. Liu, Y. Yang (2019) “A 1.2-V 2.41-GHz three-stage CMOS OTA with efficient frequency compensation technique”  IEEE Transactions on Circuits and Systems I: Regular Papers, 66(1), 20–30
  \bibitem{ref10} Feng, X., Yuan, X., Zhou, J., An, K., Zhu, F., Wei, X., Huang, Y., Zhang, J., Chen, L., Liu, J. and Li, C., (2025) “A review: CNT/diamond composites prepared via CVD and its potential applications” Materials Science in Semiconductor Processing, 186, p.109008.
  \bibitem{ref11} Okasha, N.M., Mirrashid, M., Naderpour, H., Ciftcioglu, A.O., Meddage, D.P.P. and Ezami, N., (2024) “Machine learning approach to predict the mechanical properties of cementitious materials containing carbon nanotubes” Developments in the Built Environment, 19, p.100494.
  \bibitem{ref12} S. Ijima, (2002) “Carbon nanotubes: Past, present, and future” Physica B: Condensed Matter, 323(1–4), 1–5
  \bibitem{ref13} Antil, B., Elkasabi, Y., Strahan, G.D. and Vander Wal, R.L., (2025) “Development of graphitic and non-graphitic carbons using different grade biopitch sources” Carbon, 232, p.119770.
  \bibitem{ref14} Y. B. Kim (2011) “Integrated circuit design based on carbon nanotube field effect transistor” Trans. Electr. Electron. Mater. 12 175–88
  \bibitem{ref15} J. Appenzeller, (2008) “Carbon nanotubes for high performance electronics” Proceedings of the IEEE 96(2) 206
  \bibitem{ref16} Javey, A., Guo, J., Farmer, D.B., Wang, Q., Yenilmez, E., Gordon, R.G., Lundstrom, M. and Dai, H., (2004) “Self-aligned ballistic molecular transistors and electrically parallel nanotube arrays” Nano letters, 4(7), pp.1319-1322.
  \bibitem{ref17} J. Deng, H. S. P. Wong (2007) “A compact SPICE model for carbon-nanotube field-effect transistors including nonidealities and its application—Part I: Model of the intrinsic channel region” IEEE Transactions on Electron Devices, 54(12), 3186– 3194
  \bibitem{ref18} J. Deng, H. S. P. Wong (2007) “A compact SPICE model for carbon-nanotube field-effect transistors including nonidealities and its application—Part II: Full device model and circuit performance benchmarking” IEEE Transactions on Electron Devices”,54(12), 3195–3205
  \bibitem{ref19} Marani, R. and Perri, A.G., (2024) “Design of CNTFET-Based Electronic Circuits: A Review. International Journal of Nanoscience, 23(01), p.2330007.
  \bibitem{ref20} Deng, J., Lin, A., Wan, G.C. and Wong, H.S.P., (2008) “Carbon nanotube transistor compact model for circuit design and performance optimization” ACM Journal on Emerging Technologies in Computing Systems (JETC), 4(2), pp.1-20.
  \bibitem{ref21} H. A. Javey, H. Kim, M. Brink, Q. Wang, A. Ural, J. Guo, P. McIntyre, P. McEuen, M. Lundstrom and H. Dai (2002) “High-kappa dielectrics for advanced carbon-nanotube transistors and logic gates” Nat. Mater. 1, 241-246
  \bibitem{ref22} A. Corletto, J.G. Shapter (2020) “Nanoscale Patterning of Carbon Nanotubes: Techniques, Applications, and Future” Adv Sci (Weinh), Nov 23;8(1):2001778
  \bibitem{ref23} G. E. Moore (1998) “Cramming More Components Onto Integrated Circuits” Proc. IEEE 1998, 86, 82-85
  \bibitem{ref24} J. Zou, Q. Zhang (2021) “Advances and Frontiers in Single-Walled Carbon Nanotube Electronics” Wiley, Adv Sci (Weinh),  8(23):e2102860
  \bibitem{ref25} M. Cen, S. Song, C. Cai (2017) “A high performance CNFET-based operational transconductance amplifier and its applications” Analog Integr Circ Sig Process, 9, 463-472
  \bibitem{ref26} S. M. A. Zanjani, M. Dousti, M. Dolatshahi (2019) “A new low-power, universal, multi-mode Gm-C filter in CNTFET technology” Microelectronics Journal, 90, 342-352
  \bibitem{ref27} D. S. Bethune, C. H. Kiang, M.S. Devries, G. Gorman, R. Savoy, J. Vazquez, R. Beyers, (1993) “Cobalt catalysed growth of carbon nanotube with single atomic layer walls”, Nature, Vol. 363, pp. 605-607
  \bibitem{ref28} Hashmi, F., Nizamuddin, M. and Zaidi, A., 2022, November. 9 V Graphene Transistor Based Triple Cascode Operational Transconductance Amplifier: A Comparative Analysis. In International Conference on Nanotechnology: Opportunities and Challenges (pp. 249-253). Singapore: Springer Nature Singapore.
  \bibitem{ref29} Arden, W.M., (2002) “The international technology roadmap for semiconductors—perspectives and challenges for the next 15 years” Current Opinion in Solid State and Materials Science, 6(5), pp.371-377.
  \bibitem{ref30} A Quick User Guide on Stanford University Carbon Nanotube Field Effect Transistors (CNFET) HSPICE Model V. 2.2.1.
  \bibitem{ref31} A. Balijepalli, S. Sinha, Y. Cao (2007) “Compact modeling of carbon nanotube transistor for early stage process-design” exploration Proc. of Int. Symp. on Low Power Electronics and Design (Portland, OR) pp 2–7
  \bibitem{ref32} Imran, A., Hasan, M., Islam, A. and Abbasi, S.A., (2012) “Optimized design of a 32-nm CNFET-based low-power ultrawideband CCII” IEEE transactions on Nanotechnology, 11(6), pp.1100-1109.
  \bibitem{ref33} Hashmi, F., Nizamuddin, M. and Zaidi, A., (2022) “. 9 V Graphene Transistor Based Triple Cascode Operational Transconductance Amplifier: A Comparative Analysis” In International Conference on Nanotechnology: Opportunities and Challenges (pp. 249-253). Singapore: Springer Nature Singapore.
  \bibitem{ref34} Y. Sun, V.  Kursun (2011) “N-type carbon-nanotube MOSFET device profile optimization for very large scale integration” Trans. Electr. Electron. Mater. 1243–50
  \bibitem{ref35} Paul L. McEuen, Michael Fuhrer, Hongkun Park, (2002) “Single-Walled Carbon Nanotube Electronics”, IEEE Transactions on Nanotechnology 1 (1)
  \bibitem{ref36} M. Nizamuddin, S. A. Loan, A. R. M. Alamoud and A. G. Alharbi (2017) “Design, simulation and the comparative analysis of CNTFET based multistage operational amplifiers” Journal of Nanoelectronics and Optoelectronics 12,1-11
  \bibitem{ref37} Ali, Z., Nizamuddin, M., Prasad, D., Ravi, A. and Hashmi, F., (2025) “Carbon Nanotube FET Based High Performance Folded Cascode Op-Amp with Gain Enhancement-Performance Evaluation, Comparison and Optimization” Journal of Nanoelectronics and Optoelectronics, 20(7), pp.756-766.
  \bibitem{ref38} Hashmi, F., Nizamuddin, M., Farshori, M.A., Amin, S.U. and Khan, Z.I., (2024) “Graphene nanoribbon FET technology‐based OTA for optimizing fast and energy‐efficient electronics for IoT application: Next‐generation circuit design” Micro \& Nano Letters, 19(6), p.e70002.
  \bibitem{ref39} Fuad, M.H., Nayan, M.F., Yeassin, R., Raihan, M.A., Noor, S.S. and Mahmud, R.R., 2024, March. Quantum insights into dielectric materials and oxide thickness-dependent conductance in single-walled CNTFET: a parametric simulation study. In 2024 International Conference on Advances in Computing, Communication, Electrical, and Smart Systems (iCACCESS) (pp. 01-06). IEEE.
  \bibitem{ref40} Yasir, M. and Alam, N., 2020. Systematic design of CNTFET based OTA and Op amp using gm/ID technique. Analog Integrated Circuits and Signal Processing, 102(2), pp.293-307
  \bibitem{ref41} Zhang, X., Maundy, B.J., El-Masry, E.I. and Finvers, I.G., 2000, May. A novel low-voltage operational transconductance amplifier and its applications. In 2000 IEEE International Symposium on Circuits and Systems (ISCAS) (Vol. 2, pp. 661-664). IEEE.
  \bibitem{ref42} Sayed, S.I., Abutaleb, M.M. and Nossair, Z.B., 2016. Optimization of CNFET parameters for high performance digital circuits. Advances in Materials Science and Engineering, 2016(1), p.6303725.
  \bibitem{ref43} Song, J., Kim, D.H., Tiepelt, J., Jo, Y.M., McGrath, G., Song, M., Chen, T., Wang, J., Coto, A., Palani, S. and Koh, D., 2026. Tunable and highly sensitive functionalized carbon-nanotube-based integrated systems for chemical gas sensing. Nature Sensors, pp.1-9.
  \bibitem{ref44} Kouma, J.N., Mbey, C.F., Libouga, I.O., Tchakounte, A. and Boum, A.T., 2025. Optimization of Ballistic Carbon Nanotube Field-Effect Transistor Using Response Surface Methodology for Enhanced Current Ratio Performance. ECS Journal of Solid State Science and Technology, 14(3), p.031008.
  \bibitem{ref45} Islam, M.B., Nizamuddin, M. and Siddiqui, M.S., 2025. Three Stage Class-A Operational Amplifier based on Reverse Nested Indirect Feedback Frequency Compensation using CNFETs. Micro and Nanostructures, p.208453.
\end{thebibliography}